% GENERATED FILE. Edit canonical lessons, then run npm run generate:books.
% canonical-source-hash: fnv1a64:7cb11a2abbfec761
% canonical-lessons: GE-C431-beschwerde, GE-C431-beschweren, GE-C431-reklamieren, GE-C431-meckern, GE-C431-unzufrieden

\chapter{Nicht zufrieden}
\label{ch:ge-nicht-zufrieden}
\hlchaptermodality{431}

\begin{chapteropening}
\textbf{By the end of this chapter:} \emph{I can complain: the complaint, complaining, taking faulty goods back, grumbling and being dissatisfied.}

Complete the last lesson of chapter 431: I can complain: the complaint, complaining, taking faulty goods back, grumbling and being dissatisfied.
\end{chapteropening}

\section[die Beschwerde]{die Beschwerde — the complaint}

\label{lesson:GE-C431-beschwerde}



\begin{quote}
Before the new one: say the German for to bite, then the German for to crawl.
\end{quote}

\subsection*{You'll want to know: die Beschwerde}
\textbf{die Beschwerde} — \textquotedblleft{}the complaint\textquotedblright{}.

\begin{quote}
Der Gast schreibt eine Beschwerde an das Hotel.
\end{quote}

\begin{quote}
\emph{The guest writes a complaint to the hotel.}
\end{quote}

Plural: \emph{die Beschwerden}. \emph{Beschwerden} also means aches, symptoms: \emph{Magenbeschwerden}.

\begin{cousinweb}
\emph{Beschwerde} is built on \emph{beschweren} 'to weigh down, to burden', from \emph{schwer} 'heavy', inherited from Old High German \emph{swār} (EtymDB): a complaint is a weight put on someone.
\end{cousinweb}

\begin{grammarlens}[title={What you've built}]
One more word for telling, explaining and describing at B1.
\end{grammarlens}

\subsection*{Your turn}
\begin{itemize}
\raggedright
  \item \emph{Say it:} \emph{die Beschwerde}
  \item \emph{Say it:} \emph{die Beschwerde}, once more
  \item \emph{Say it:} a sentence of your own with \emph{die Beschwerde}
  \item \emph{Recall:} say the German for to bite, then the German for to crawl, then say \emph{die Beschwerde} again
\end{itemize}

\begin{tcolorbox}[breakable,colback=teal!4,colframe=teal!35!black,title={Before you move on}]
What does die Beschwerde mean? (\textquotedblleft{}The complaint\textquotedblright{}.) Say it once more.
\end{tcolorbox}
\section[beschweren]{beschweren — to complain (mostly reflexive: sich beschweren)}

\label{lesson:GE-C431-beschweren}



\begin{quote}
Before the new one: say the German for to crawl, then the German for the complaint.
\end{quote}

\subsection*{You'll want to know: beschweren}
\textbf{beschweren} — \textquotedblleft{}to complain (mostly reflexive: sich beschweren)\textquotedblright{}.

\begin{quote}
Sie beschwert sich über den Lärm im Flur.
\end{quote}

\begin{quote}
\emph{She complains about the noise in the corridor.}
\end{quote}

Reflexive, with \emph{über}: \emph{ich beschwere mich über das Essen}. Without \emph{sich} it means to weigh down.

\begin{cousinweb}
\emph{beschweren} is built from \emph{be-} and \emph{schwer} 'heavy', inherited from Old High German \emph{swār} (EtymDB); the reflexive sense 'to complain' grew from 'to feel weighed down'.
\end{cousinweb}

\begin{grammarlens}[title={What you've built}]
One more word for telling, explaining and describing at B1.
\end{grammarlens}

\subsection*{Your turn}
\begin{itemize}
\raggedright
  \item \emph{Say it:} \emph{beschweren}
  \item \emph{Say it:} \emph{beschweren}, once more
  \item \emph{Say it:} a sentence of your own with \emph{beschweren}
  \item \emph{Recall:} say the German for to crawl, then the German for the complaint, then say \emph{beschweren} again
\end{itemize}

\begin{tcolorbox}[breakable,colback=teal!4,colframe=teal!35!black,title={Before you move on}]
What does beschweren mean? (\textquotedblleft{}To complain (mostly reflexive: sich beschweren)\textquotedblright{}.) Say it once more.
\end{tcolorbox}
\section[reklamieren]{reklamieren — to take back (faulty goods), to make a complaint about}

\label{lesson:GE-C431-reklamieren}



\begin{quote}
Before the new one: say the German for the complaint, then the German for to complain.
\end{quote}

\subsection*{You'll want to know: reklamieren}
\textbf{reklamieren} — \textquotedblleft{}to take back (faulty goods), to make a complaint about\textquotedblright{}.

\begin{quote}
Er reklamiert die neue Kamera, weil sie nicht funktioniert.
\end{quote}

\begin{quote}
\emph{He takes the new camera back because it does not work.}
\end{quote}

Past participle without \emph{ge-}: \emph{ich habe reklamiert}. \emph{Die Reklamation} — the complaint about goods.

\begin{cousinweb}
\emph{reklamieren} was borrowed from French \emph{réclamer} (EtymDB), which goes back to Latin \emph{reclāmāre} 'to cry out against, to protest', from \emph{clāmāre} 'to shout' (Lewis \& Short).
\end{cousinweb}

\begin{grammarlens}[title={What you've built}]
One more word for telling, explaining and describing at B1.
\end{grammarlens}

\subsection*{Your turn}
\begin{itemize}
\raggedright
  \item \emph{Say it:} \emph{reklamieren}
  \item \emph{Say it:} \emph{reklamieren}, once more
  \item \emph{Say it:} a sentence of your own with \emph{reklamieren}
  \item \emph{Recall:} say the German for the complaint, then the German for to complain, then say \emph{reklamieren} again
\end{itemize}

\begin{tcolorbox}[breakable,colback=teal!4,colframe=teal!35!black,title={Before you move on}]
What does reklamieren mean? (\textquotedblleft{}To take back (faulty goods), to make a complaint about\textquotedblright{}.) Say it once more.
\end{tcolorbox}
\section[meckern]{meckern — to grumble, to moan}

\label{lesson:GE-C431-meckern}



\begin{quote}
Before the new one: say the German for to complain, then the German for to take back.
\end{quote}

\subsection*{You'll want to know: meckern}
\textbf{meckern} — \textquotedblleft{}to grumble, to moan\textquotedblright{}.

\begin{quote}
Mein Onkel meckert oft über das Wetter.
\end{quote}

\begin{quote}
\emph{My uncle often moans about the weather.}
\end{quote}

Everyday speech. First the sound a goat makes.

\begin{cousinweb}
\emph{meckern} first described the bleating of a goat and is usually explained as a sound word. EtymDB lists it without a source.
\end{cousinweb}

\begin{grammarlens}[title={What you've built}]
One more word for telling, explaining and describing at B1.
\end{grammarlens}

\subsection*{Your turn}
\begin{itemize}
\raggedright
  \item \emph{Say it:} \emph{meckern}
  \item \emph{Say it:} \emph{meckern}, once more
  \item \emph{Say it:} a sentence of your own with \emph{meckern}
  \item \emph{Recall:} say the German for to complain, then the German for to take back, then say \emph{meckern} again
\end{itemize}

\begin{tcolorbox}[breakable,colback=teal!4,colframe=teal!35!black,title={Before you move on}]
What does meckern mean? (\textquotedblleft{}To grumble, to moan\textquotedblright{}.) Say it once more.
\end{tcolorbox}
\section[unzufrieden]{unzufrieden — dissatisfied, unhappy (with something)}

\label{lesson:GE-C431-unzufrieden}



\begin{quote}
Before the new one: say the German for to take back, then the German for to grumble.
\end{quote}

\subsection*{You'll want to know: unzufrieden}
\textbf{unzufrieden} — \textquotedblleft{}dissatisfied, unhappy (with something)\textquotedblright{}.

\begin{quote}
Viele Gäste sind mit dem Frühstück unzufrieden.
\end{quote}

\begin{quote}
\emph{Many guests are unhappy with the breakfast.}
\end{quote}

With \emph{mit}: \emph{unzufrieden mit dem Service}.

\begin{cousinweb}
\emph{unzufrieden} is built from \emph{un-} 'not' and \emph{zufrieden} 'content' (EtymDB), which comes from \emph{zu Frieden} 'at peace'.
\end{cousinweb}

\begin{grammarlens}[title={What you've built}]
One more word for telling, explaining and describing at B1.
\end{grammarlens}

\subsection*{Your turn}
\begin{itemize}
\raggedright
  \item \emph{Say it:} \emph{unzufrieden}
  \item \emph{Say it:} \emph{unzufrieden}, once more
  \item \emph{Say it:} a sentence of your own with \emph{unzufrieden}
  \item \emph{Recall:} say the German for to take back, then the German for to grumble, then say \emph{unzufrieden} again
\end{itemize}

\begin{tcolorbox}[breakable,colback=teal!4,colframe=teal!35!black,title={Before you move on}]
What does unzufrieden mean? (\textquotedblleft{}Dissatisfied, unhappy (with something)\textquotedblright{}.) Say it once more.
\end{tcolorbox}
